% TOP_PROBLEM: {"id": 1187, "title": "The Andrews-Curtis Conjecture", "category_id": 4, "category": {"id": 4, "name": "algebra", "display_name": "Algebra", "description": "Group theory, ring theory, field theory, and algebraic structures.", "slug": "algebra", "order_index": 4, "created_at": "2026-07-31T15:26:25.671Z"}, "status": "open"}
% FIRST_POSTED: null
% ENTRY_KIND: statement_only
% Imported from: batch04.tex; UnsolvedMath ID 1187
\documentclass[11pt]{article}
\usepackage[margin=25mm]{geometry}
\usepackage{fontspec}
\setmainfont{Latin Modern Roman}
\usepackage{amsmath,amssymb,mathtools}
\usepackage{unicode-math}
\setmathfont{Latin Modern Math}
\usepackage{xurl,hyperref}
\hypersetup{colorlinks=true,urlcolor=blue}
\setcounter{secnumdepth}{0}
\setlength{\parindent}{0pt}
\setlength{\parskip}{5pt}
\setlength{\emergencystretch}{3em}
\newcommand{\sref}[2]{\hyperlink{source:#1:#2}{[S#2]}}
\newcommand{\cataloguescope}[1]{\paragraph{Recorded target.}#1}
\begin{document}
% UnsolvedMath ID 1187; source catalogue code ALG-023
\setcounter{section}{1186}
\section{The Andrews-Curtis Conjecture}\label{1187}
{\small\sffamily UnsolvedMath ID 1187 · Algebra}
\subsection{Definitions and mathematical statement}

\paragraph{Shared notation.}$\mathbb N=\{1,2,\ldots\}$, and $\mathbb Z,\mathbb Q,\mathbb R,\mathbb C$ denote the integers, rationals, reals and complex numbers. Any other conventions are specified locally.


Let \(F_r=F(x_1,\ldots,x_r)\) be the free group on \(r\ge1\)
generators.  A balanced presentation of the trivial group is
\[
 \langle x_1,\ldots,x_r\mid w_1,\ldots,w_r\rangle=\{1\},
\]
or equivalently the normal closure of \(w_1,\ldots,w_r\) is all of
\(F_r\).  An elementary Andrews--Curtis move on the relator tuple
replaces one \(w_i\), leaving every other entry fixed, by one of
\[
 w_iw_j\ (j\ne i),\qquad w_i^{-1},\qquad
 v w_i v^{-1}\ (v\in F_r).
\]
Words represent elements of the free group, so free reductions are implicit.
The standard tuple is \((x_1,\ldots,x_r)\).
The number of generators and relators remains fixed during the moves.
\paragraph{Statement recorded in the supplied source.}
For every \(r\ge1\) and every balanced presentation of the trivial
group on generators \(x_1,\ldots,x_r\), can its relator tuple
\((w_1,\ldots,w_r)\) be transformed into
\((x_1,\ldots,x_r)\) by finitely many elementary
Andrews--Curtis moves? \sref{1187}{2}
\subsection{Short English statement}
Can every balanced presentation of the trivial group be reduced to the standard presentation using relator multiplication, inversion and conjugation, without adding generators?
\subsection{Sources}
\begin{description}
\item[\hypertarget{source:1187:1}{S1}] ULAM.AI and the UnsolvedMath Contributors, \emph{UnsolvedMath: A Curated Collection of Open Mathematics Problems}, record 1187 (ALG-023). CC BY 4.0. \url{https://huggingface.co/datasets/ulamai/UnsolvedMath}.
\item[\hypertarget{source:1187:2}{S2}] Alexei D. Miasnikov and Alexei G. Myasnikov, \emph{Balanced presentations of the trivial group on two generators and the Andrews--Curtis conjecture}, Groups and Computation III, volume 23 (2001), 257--263, Section 1, transformations (AC1)--(AC3) and the conjecture; arXiv:math/0304305. \url{https://arxiv.org/abs/math/0304305}.
\end{description}
\label{end:1187}
\end{document}
